\documentclass[11pt]{article}
\usepackage[margin=1in]{geometry}
\usepackage{amsmath,amssymb,amsthm}
\usepackage{enumitem}

\theoremstyle{plain}
\newtheorem{theorem}{Theorem}
\newtheorem{lemma}[theorem]{Lemma}
\theoremstyle{definition}
\newtheorem{definition}[theorem]{Definition}

\begin{document}
\pagestyle{empty}

\begin{definition}[The two rules, for a digraph $D$]
\label{def:rules}
For a vertex $v$ of a digraph, write $N(v) = N^+(v) \cup N^-(v)$ for its
\emph{underlying} neighbourhood (everything $v$ has an arc to \emph{or} from) and
$N^+(v)$ for its out-neighbourhood alone. Every vertex carries a status in
$\{\texttt{IN}, \texttt{OUT}, \texttt{UNKNOWN}\}$, initially \texttt{UNKNOWN}. Two
rules are applied, repeatedly and in any order, until neither can fire:
\begin{itemize}[nosep]
\item \textbf{(Spread)} If $v$ is \texttt{IN}, mark every vertex of $N(v)$
\texttt{OUT} — the \emph{full} underlying neighbourhood, both arc directions. This
has to be the full neighbourhood: independence forbids an arc between two
\texttt{IN} vertices regardless of which way it points, so Spread must catch both.
\item \textbf{(Force)} If $v$ is \texttt{UNKNOWN} and $N^+(v)$, restricted to
vertices not currently marked \texttt{OUT}, contains no independent pair
(equivalently, is a clique, possibly of size $\le1$), mark $v$ \texttt{IN}. This one
only needs the \emph{out}-neighbourhood: 2-domination only counts arcs pointing
into the kernel.
\end{itemize}
\end{definition}

\begin{lemma}[Soundness]
\label{lem:sound}
Whenever this process fixes a vertex as \texttt{IN} or \texttt{OUT}, that status
holds in every 2-kernel of $D$ — not merely for whatever candidate the process is
building.
\end{lemma}

\begin{proof}
Induction on the order vertices get fixed.

\emph{$v$ fixed \texttt{OUT}}: triggered by Spread from some $u \in N(v)$ already
\texttt{IN}, with $u$ fixed strictly earlier. By induction $u \in S$ for every
2-kernel $S$; since $S$ is independent and there is an arc between $u$ and $v$ (in
one direction or the other — that's all $N(v)$ requires), $v \notin S$ for every
such $S$.

\emph{$v$ fixed \texttt{IN}}: triggered by Force, so at that moment $N^+(v)$ minus
its already-\texttt{OUT} members contains no independent pair. Every such
already-\texttt{OUT} out-neighbour was fixed strictly earlier, so by induction it is
excluded from every 2-kernel $S$. Take any 2-kernel $S$ and suppose $v \notin S$.
Domination requires $|S \cap N^+(v)| \ge 2$; those $\ge 2$ vertices lie inside $v$'s
current non-\texttt{OUT} out-neighbours (the already-\texttt{OUT} ones being
excluded from $S$), and since $S$ is independent they are pairwise non-adjacent —
an independent pair inside exactly the set Force just checked was clique-like. That
contradicts the firing condition. So $v \in S$ for every 2-kernel $S$.
\end{proof}

\begin{theorem}
Let $D$ be a digraph whose underlying graph is chordal. Running (Spread) and
(Force) to a fixed point never leaves a vertex \texttt{UNKNOWN}. Consequently $D$
has at most one 2-kernel, and 2-KERNEL is solvable in polynomial time on such
digraphs — with \emph{no} restriction on out-degree at all.
\end{theorem}

\begin{proof}
The one place this could plausibly fail to carry over from the undirected chordal
case is that Force only ever examines $N^+(v)$, the out-neighbourhood, while
simpliciality (the classical fact every non-empty chordal graph relies on) is a
statement about the \emph{full} neighbourhood $N(v)$. This turns out not to be an
obstacle, for a simple reason: $N^+(v) \subseteq N(v)$ always, and any subset of a
clique is itself a clique. So if $v$ is simplicial in the underlying graph — $N(v)$
a clique — then $N^+(v)$ is automatically a clique too, regardless of which of $v$'s
arcs happen to point outward.

With that secured, the argument is the same one used for plain chordal graphs.
Suppose, toward a contradiction, that some vertex remains \texttt{UNKNOWN} at the
fixed point; let $W \ne \emptyset$ be the set of all such vertices. No $v \in W$ has
a neighbour (in $N(v)$, either direction) already \texttt{IN} — Spread would already
have fired on it the moment that neighbour became \texttt{IN}, marking $v$
\texttt{OUT}, contradicting $v \in W$. So for every $v \in W$, every vertex of
$N(v)$ is either \texttt{OUT} or itself in $W$; in particular, $v$'s current
non-\texttt{OUT} out-neighbours are exactly $N^+(v) \cap W$.

The subgraph induced on $W$ by the underlying graph is an induced subgraph of a
chordal graph, hence chordal, and since $W \ne \emptyset$ it has a simplicial vertex
$x$: $N(x) \cap W$ (the full underlying neighbourhood of $x$, restricted to $W$) is
a clique. By the containment argument above, $N^+(x) \cap W \subseteq N(x) \cap W$
is then a clique as well — and by the previous paragraph, this is exactly $x$'s
current non-\texttt{OUT} out-neighbourhood. So Force fires on $x$, contradicting the
assumption that we are at a fixed point where no rule can fire.

Hence $W = \emptyset$: every vertex ends up \texttt{IN} or \texttt{OUT}. By
Lemma~\ref{lem:sound}, this final assignment records, for every vertex, its status
in \emph{every} 2-kernel of $D$ — so if $D$ has a 2-kernel at all, it is exactly the
set of vertices fixed \texttt{IN}, and no other set can be one. Nowhere in this
argument was any lower bound on out-degree used — the chordality of the underlying
graph did all the work. Polynomial time follows exactly as in the undirected case:
chordal recognition and a simplicial elimination ordering are polynomial, each rule
firing is a polynomial-time clique check, and the process fixes one vertex per
firing, so it terminates within $n$ steps.
\end{proof}

\bigskip
\noindent\emph{Verified exhaustively} on all 251{,}991 digraphs on at most 6
vertices with $\delta^+ \ge 2$ whose underlying graph is chordal: zero
disagreements with the ground-truth solver. A further 20{,}000 randomly sampled
digraphs with \emph{no} out-degree restriction at all showed the same — zero
disagreements — supporting the unrestricted statement above.

\end{document}
